% Source: https://artofproblemsolving.com/community/c3h2433694_forester2015s_problems
% Stable identity: aops-topic-2433694-q39
\textbf{Upper bound.} Two distinct lines meet in at most one point, so $N\le\dbinom42=6$.

\textbf{Possible values.} The configurations below give $N=0$ (four parallel lines), $1$ (four lines through one point), $3$ (three parallel lines and a transversal), $4$ (two pairs of parallel lines), $5$ (two parallel lines and two lines crossing each other between them) and $6$ (no two parallel, no three through a point).

\begin{asy}
usepackage("amsmath");
size(7cm);
defaultpen(fontsize(10pt));
picture[] p;
for(int i=0;i<6;++i) p.push(new picture);
// N=0
for(int k=0;k<4;++k) draw(p[0],(-1,-1+0.6k)--(1,-0.4+0.6k));
label(p[0],"$N=0$",(0,-1.6));
// N=1
for(int k=0;k<4;++k) draw(p[1],1.1*dir(20+45k)--1.1*dir(200+45k)); dot(p[1],(0,0));
label(p[1],"$N=1$",(0,-1.6));
// N=3: three parallel, one transversal
for(int k=0;k<3;++k) draw(p[2],(-1,-0.9+0.7k)--(1,-0.9+0.7k)); draw(p[2],(-0.6,-1.2)--(0.6,1.0));
label(p[2],"$N=3$",(0,-1.6));
// N=4: two pairs of parallel lines
draw(p[3],(-1,-0.4)--(1,-0.4)); draw(p[3],(-1,0.4)--(1,0.4)); draw(p[3],(-0.7,-1)--(-0.2,1)); draw(p[3],(0.2,-1)--(0.7,1));
label(p[3],"$N=4$",(0,-1.6));
// N=5: two parallel, two others crossing between them
draw(p[4],(-1,-0.6)--(1,-0.6)); draw(p[4],(-1,0.6)--(1,0.6)); draw(p[4],(-0.8,-1)--(0.8,1)); draw(p[4],(0.8,-1)--(-0.8,1));
label(p[4],"$N=5$",(0,-1.6));
// N=6: general position
draw(p[5],(-1,-0.8)--(1,0.2)); draw(p[5],(-1,0.5)--(1,-0.9)); draw(p[5],(-0.7,-1.1)--(0.2,1.1)); draw(p[5],(0.9,-1.1)--(0.1,1.1));
label(p[5],"$N=6$",(0,-1.6));
for(int i=0;i<6;++i) add(shift(2.6*(i%3),-3.2*floor(i/3))*p[i]);
\end{asy}

\textbf{$N=2$ is impossible.} Suppose the only intersection points are $P$ and $Q$.
\begin{itemize}
\item \emph{Some line $\ell$ passes through both.} At least one other line $m$ passes through $P$ and at least one other line $n$ passes through $Q$; $m\ne n$, since otherwise this line would pass through $P$ and $Q$ and equal $\ell$. If $m$ and $n$ met, they would meet at $P$ or $Q$, and again one of them would equal $\ell$; so $m\parallel n$. The fourth line $k$ is not parallel to $m$: otherwise it would cross $\ell$ at $P$ or $Q$ and coincide with $m$ or $n$. So $k$ meets $m$ and $n$, necessarily at $P$ and at $Q$, and then $k=\ell$, a contradiction.
\item \emph{No line passes through both.} Then exactly two lines pass through $P$ and the other two through $Q$. A line through $P$ and a line through $Q$ cannot meet (their common point would be $P$ or $Q$), so they are parallel. Then both lines through $P$ are parallel to the same line through $Q$, hence to each other, although they meet at $P$: a contradiction.
\end{itemize}
So the possible values are $0,1,3,4,5,6$, with sum $19$. The answer is $\boxed{\textbf{(D)}\ 19}$.
